
\documentclass{article}
\usepackage{colortbl}
\usepackage{makecell}
\usepackage{multirow}
\usepackage{supertabular}

\begin{document}

\newcounter{utterance}

\centering \large Interaction Transcript for game `dond', experiment `coop\_en', episode 13 with reflexion{-}Qwen3.5{-}4B.
\vspace{24pt}

{ \footnotesize  \setcounter{utterance}{1}
\setlength{\tabcolsep}{0pt}
\begin{supertabular}{c@{$\;$}|p{.15\linewidth}@{}p{.15\linewidth}p{.15\linewidth}p{.15\linewidth}p{.15\linewidth}p{.15\linewidth}}
    \# & $\;$A & \multicolumn{4}{c}{Game Master} & $\;\:$B\\
    \hline

    \theutterance \stepcounter{utterance}  
    & & \multicolumn{4}{p{0.6\linewidth}}{
        \cellcolor[rgb]{0.9,0.9,0.9}{
            \makecell[{{p{\linewidth}}}]{
                \texttt{\tiny{[P1$\langle$GM]}}
                \texttt{You have attempted this same task before and did not succeed. These are the plans you wrote after those attempts.} \\
\texttt{Trial \#1: New plan: I failed because I assumed the other player would accept my offer of combs for a statue and a stamp, but I did not verify if their secret valuation of the combs was high enough to make that trade mutually beneficial. Since the goal is to maximize the *sum* of both scores, I should have calculated the total value of all items first (statue: 1, combs: 8, stamp: 1, total: 10) and then proposed a split that divides this total as evenly as possible (e.g., 5 each), rather than focusing solely on my own high{-}value items. Next time, I will first calculate the total value of the set, then propose a division where the sum of my items plus the sum of the other player's items equals the total, ensuring the split is roughly 50/50 or the most equitable distribution possible, regardless of which specific items I hold.} \\
\\ 
\texttt{Trial \#2: New plan: I failed because I proposed a split (2 combs, 1 statue) that left the other player with only 1 stamp and 2 screws. Even if the other player values combs highly, they received items I valued at 0 (screws) and low value (stamp), likely resulting in a total score for them that was too low to make the deal mutually beneficial or complementary. The goal is to maximize the *sum* of both scores, which means finding a split where the other player receives enough value to accept. I should have calculated the total value of all items (2 statues, 4 combs, 1 stamp, 2 screws = 20 total value assuming symmetric values, or at least ensuring the other player gets a significant portion). Next time, I will propose a split where the other player receives the items I value least (screws) plus one or two other items (like the stamp or one comb) to ensure their total value is high enough to make the agreement successful, rather than giving them only the lowest{-}value items. Specifically, I will propose: {[}Proposal: 2 comb, 1 screw{]} for myself, leaving the other player with 1 statue, 1 comb, 1 stamp, and 1 screw (or similar balanced distribution) to ensure the sum is maximized and the deal is accepted.} \\
\\ 
\texttt{Trial \#3: New plan: I failed because I included square brackets in my message, which the Game Master interpreted as a secret proposal rather than a negotiation message, causing the message to be rejected. The rules state that messages to the other player must not use square brackets. Next time, I will communicate my offer directly without brackets, clearly stating the items I want to keep and the items I am offering to the other player, ensuring the format matches the "negotiation" phase and not the "secret proposal" phase. For example, I will say: "I propose I keep the statue, 1 comb, and 2 screws. You take the remaining 2 combs and 1 stamp." This avoids the syntax error while still conveying the necessary information for a potential agreement.} \\
\\ 
\\ 
\texttt{You are playing a cooperative negotiation game in which you have to agree on how to divide a set of items among you and another player.} \\
\\ 
\texttt{Rules:} \\
\texttt{(a) You and the other player are given a set of items. Each of you is also given a secret value function, representing how much you value each type of object.} \\
\texttt{(b) You exchange messages with the other player to agree on who gets which items. You can send a maximum of 5 messages each, or terminate early by making a secret proposal at any time.} \\
\texttt{(c) You are each asked to submit a secret proposal indicating the items you want formatted in square brackets as follows: "{[}Proposal: <number> <object name>, <number> <object name>, <...>{]}"} \\
\texttt{(d) If your proposals are complementary, i.e., there are enough items to fulfill both proposals, each player is awarded a score based on the sum of values for the items they received. Otherwise, both of you get zero points.} \\
\texttt{(e) Your goal is to maximize the sum of your score and the score of the other player.} \\
\\ 
\texttt{Let us start.} \\
\\ 
\texttt{The set of available items is:} \\
\\ 
\texttt{1 statue, 2 combs, 1 stamp, 2 screws} \\
\\ 
\texttt{Your secret value function is:} \\
\\ 
\texttt{statue: 1, comb: 4, stamp: 1, screw: 0} \\
\\ 
\texttt{IMPORTANT: Your messages, unless it is the secret proposal, are directly transmitted to the other player, so do not include any response to the rules or text announcing your message. To make a secret proposal, use the indicated format. Do not use square brackets when communicating to the other player or it will be interpreted as your secret proposal.} \\
            }
        }
    }
    & \\ \\

    \theutterance \stepcounter{utterance}  
    & \multicolumn{4}{p{0.6\linewidth}}{
        \cellcolor[rgb]{0.9,0.9,0.9}{
            \makecell[{{p{\linewidth}}}]{
                \texttt{\tiny{[P1$\rangle$GM]}}
                \texttt{I propose I keep the statue, 2 combs, and 1 stamp. You take the 2 screws.} \\
            }
        }
    }
    & & \\ \\

    \theutterance \stepcounter{utterance}  
    & & & \multicolumn{2}{p{0.3\linewidth}}{
        \cellcolor[rgb]{0.9,0.9,0.9}{
            \makecell[{{p{\linewidth}}}]{
                \texttt{\tiny{[GM$|$GM]}}
                \texttt{continue} \\
            }
        }
    }
    & & \\ \\

    \theutterance \stepcounter{utterance}  
    & & & \multicolumn{2}{p{0.3\linewidth}}{
        \cellcolor[rgb]{0.9,0.9,0.9}{
            \makecell[{{p{\linewidth}}}]{
                \texttt{\tiny{[GM$|$GM]}}
                \texttt{I propose I keep the statue, 2 combs, and 1 stamp. You take the 2 screws.} \\
            }
        }
    }
    & & \\ \\

    \theutterance \stepcounter{utterance}  
    & & \multicolumn{4}{p{0.6\linewidth}}{
        \cellcolor[rgb]{0.9,0.9,0.9}{
            \makecell[{{p{\linewidth}}}]{
                \texttt{\tiny{[GM$\rangle$P2]}}
                \texttt{You have attempted this same task before and did not succeed. These are the plans you wrote after those attempts.} \\
\texttt{Trial \#1: New plan: I failed because I submitted a proposal containing all items, which is not complementary to any possible proposal from the other player (since they must also receive items to get a non{-}zero score). To succeed, I must propose a subset of items that leaves enough items for the other player to form a valid, complementary set. Given my values (statue: 4, stamp: 2, combs: 1, screws: 1), I should prioritize the statue and stamp. I will propose taking the statue and one comb, leaving the stamp, remaining comb, and two screws for the other player. This ensures both players have items, maximizing the chance of a successful agreement while securing my highest{-}value items.} \\
\\ 
\texttt{Trial \#2: New plan: I failed because I proposed taking all the items I valued (statue, stamp, and screws), leaving the other player with nothing. Since the other player also needs items to receive a non{-}zero score, my proposal was not complementary to any valid response from them. To succeed, I must propose a split where I take a subset of items and leave a complementary subset for the other player. Given my values (statue: 4, stamp: 2, combs: 1, screws: 1), I should prioritize the statue and stamp. I will propose taking the statue and one comb (my value: 5), leaving the stamp, one comb, and two screws for the other player. This ensures both players have items, maximizing the chance of a successful agreement while securing my highest{-}value items.} \\
\\ 
\\ 
\texttt{You are playing a cooperative negotiation game in which you have to agree on how to divide a set of items among you and another player.} \\
\\ 
\texttt{Rules:} \\
\texttt{(a) You and the other player are given a set of items. Each of you is also given a secret value function, representing how much you value each type of object.} \\
\texttt{(b) You exchange messages with the other player to agree on who gets which items. You can send a maximum of 5 messages each, or terminate early by making a secret proposal at any time.} \\
\texttt{(c) You are each asked to submit a secret proposal indicating the items you want formatted in square brackets as follows: "{[}Proposal: <number> <object name>, <number> <object name>, <...>{]}"} \\
\texttt{(d) If your proposals are complementary, i.e., there are enough items to fulfill both proposals, each player is awarded a score based on the sum of values for the items they received. Otherwise, both of you get zero points.} \\
\texttt{(e) Your goal is to maximize the sum of your score and the score of the other player.} \\
\\ 
\texttt{Let us start.} \\
\\ 
\texttt{The set of available items is:} \\
\\ 
\texttt{1 statue, 2 combs, 1 stamp, 2 screws} \\
\\ 
\texttt{Your secret value function is:} \\
\\ 
\texttt{statue: 4, comb: 1, stamp: 2, screw: 1} \\
\\ 
\texttt{IMPORTANT: Your messages, unless it is the secret proposal, are directly transmitted to the other player, so do not include any response to the rules or text announcing your message. To make a secret proposal, use the indicated format. Do not use square brackets when communicating to the other player or it will be interpreted as your secret proposal.} \\
\\ 
\\ 
\texttt{I propose I keep the statue, 2 combs, and 1 stamp. You take the 2 screws.} \\
            }
        }
    }
    & \\ \\

    \theutterance \stepcounter{utterance}  
    & & & \multicolumn{4}{p{0.6\linewidth}}{
        \cellcolor[rgb]{0.9,0.9,0.9}{
            \makecell[{{p{\linewidth}}}]{
                \texttt{\tiny{[GM$\langle$P2]}}
                \texttt{I propose I keep the statue, 1 comb, and 1 stamp. You take the remaining comb and 2 screws.} \\
            }
        }
    }
     \\ \\

    \theutterance \stepcounter{utterance}  
    & & & \multicolumn{2}{p{0.3\linewidth}}{
        \cellcolor[rgb]{0.9,0.9,0.9}{
            \makecell[{{p{\linewidth}}}]{
                \texttt{\tiny{[GM$|$GM]}}
                \texttt{continue} \\
            }
        }
    }
    & & \\ \\

    \theutterance \stepcounter{utterance}  
    & & & \multicolumn{2}{p{0.3\linewidth}}{
        \cellcolor[rgb]{0.9,0.9,0.9}{
            \makecell[{{p{\linewidth}}}]{
                \texttt{\tiny{[GM$|$GM]}}
                \texttt{I propose I keep the statue, 1 comb, and 1 stamp. You take the remaining comb and 2 screws.} \\
            }
        }
    }
    & & \\ \\

    \theutterance \stepcounter{utterance}  
    & & \multicolumn{4}{p{0.6\linewidth}}{
        \cellcolor[rgb]{0.9,0.9,0.9}{
            \makecell[{{p{\linewidth}}}]{
                \texttt{\tiny{[P1$\langle$GM]}}
                \texttt{I propose I keep the statue, 1 comb, and 1 stamp. You take the remaining comb and 2 screws.} \\
            }
        }
    }
    & \\ \\

    \theutterance \stepcounter{utterance}  
    & \multicolumn{4}{p{0.6\linewidth}}{
        \cellcolor[rgb]{0.9,0.9,0.9}{
            \makecell[{{p{\linewidth}}}]{
                \texttt{\tiny{[P1$\rangle$GM]}}
                \texttt{I propose I keep the statue, 1 comb, and 1 stamp. You take the remaining comb and 2 screws.} \\
\\ 
\texttt{*Wait, I need to re{-}evaluate my strategy based on the rules and the goal.*} \\
\\ 
\texttt{**Analysis:**} \\
\texttt{1.  **My Values:** Statue (1), Comb (4), Stamp (1), Screw (0). Total = 6.} \\
\texttt{2.  **Other Player's Values:** Assuming symmetric values (Statue 1, Comb 4, Stamp 1, Screw 0) or at least high value for combs.} \\
\texttt{3.  **Goal:** Maximize the *sum* of both scores. This means we need a mutually beneficial deal where both players get a positive score, ideally close to the total value divided by two.} \\
\texttt{4.  **Current Situation:**} \\
\texttt{    *   I offered: Keep Statue, 1 Comb, 1 Stamp. Give them 1 Comb, 2 Screws.} \\
\texttt{    *   They offered: Keep Statue, 1 Comb, 1 Stamp. Give them 1 Comb, 2 Screws.} \\
\texttt{    *   **Conflict:** We both want the Statue and the Stamp. If we both keep the Statue, we can't both have it. If we both keep the Stamp, we can't both have it.} \\
\texttt{    *   **My Mistake:** In my previous turn, I proposed keeping the Statue, 1 Comb, and 1 Stamp. The other player also proposed keeping the Statue, 1 Comb, and 1 Stamp. This is a direct conflict over the Statue and the Stamp.} \\
\texttt{    *   **Correction:** I must yield on the high{-}value items (Statue and Stamp) to the other player if they want them, or find a split where the total value is maximized. Since I value Combs highly (4), I should try to keep the Combs. The other player likely values Combs highly too.} \\
\texttt{    *   **Optimal Split:**} \\
\texttt{        *   Total Items: 1 Statue, 2 Combs, 1 Stamp, 2 Screws.} \\
\texttt{        *   My Top Values: Combs (4 each).} \\
\texttt{        *   Other Player's Top Values: Likely Combs (4 each).} \\
\texttt{        *   If we split the 2 Combs (1 each), we get 4+4=8 value from combs alone.} \\
\texttt{        *   Statue (1) and Stamp (1) are low value (1 each).} \\
\texttt{        *   Screws (0) are worthless to me.} \\
\texttt{        *   **Best Deal:** I keep 1 Comb. Other player keeps 1 Comb. I give them the Statue and Stamp (or vice versa).} \\
\texttt{        *   Let's try: I keep 1 Comb. Other player keeps 1 Comb, Statue, Stamp.} \\
\texttt{        *   My Score: 4.} \\
            }
        }
    }
    & & \\ \\

    \theutterance \stepcounter{utterance}  
    & & & \multicolumn{2}{p{0.3\linewidth}}{
        \cellcolor[rgb]{0.9,0.9,0.9}{
            \makecell[{{p{\linewidth}}}]{
                \texttt{\tiny{[GM$|$GM]}}
                \texttt{continue} \\
            }
        }
    }
    & & \\ \\

    \theutterance \stepcounter{utterance}  
    & & & \multicolumn{2}{p{0.3\linewidth}}{
        \cellcolor[rgb]{0.9,0.9,0.9}{
            \makecell[{{p{\linewidth}}}]{
                \texttt{\tiny{[GM$|$GM]}}
                \texttt{I propose I keep the statue, 1 comb, and 1 stamp. You take the remaining comb and 2 screws.} \\
\\ 
\texttt{*Wait, I need to re{-}evaluate my strategy based on the rules and the goal.*} \\
\\ 
\texttt{**Analysis:**} \\
\texttt{1.  **My Values:** Statue (1), Comb (4), Stamp (1), Screw (0). Total = 6.} \\
\texttt{2.  **Other Player's Values:** Assuming symmetric values (Statue 1, Comb 4, Stamp 1, Screw 0) or at least high value for combs.} \\
\texttt{3.  **Goal:** Maximize the *sum* of both scores. This means we need a mutually beneficial deal where both players get a positive score, ideally close to the total value divided by two.} \\
\texttt{4.  **Current Situation:**} \\
\texttt{    *   I offered: Keep Statue, 1 Comb, 1 Stamp. Give them 1 Comb, 2 Screws.} \\
\texttt{    *   They offered: Keep Statue, 1 Comb, 1 Stamp. Give them 1 Comb, 2 Screws.} \\
\texttt{    *   **Conflict:** We both want the Statue and the Stamp. If we both keep the Statue, we can't both have it. If we both keep the Stamp, we can't both have it.} \\
\texttt{    *   **My Mistake:** In my previous turn, I proposed keeping the Statue, 1 Comb, and 1 Stamp. The other player also proposed keeping the Statue, 1 Comb, and 1 Stamp. This is a direct conflict over the Statue and the Stamp.} \\
\texttt{    *   **Correction:** I must yield on the high{-}value items (Statue and Stamp) to the other player if they want them, or find a split where the total value is maximized. Since I value Combs highly (4), I should try to keep the Combs. The other player likely values Combs highly too.} \\
\texttt{    *   **Optimal Split:**} \\
\texttt{        *   Total Items: 1 Statue, 2 Combs, 1 Stamp, 2 Screws.} \\
\texttt{        *   My Top Values: Combs (4 each).} \\
\texttt{        *   Other Player's Top Values: Likely Combs (4 each).} \\
\texttt{        *   If we split the 2 Combs (1 each), we get 4+4=8 value from combs alone.} \\
\texttt{        *   Statue (1) and Stamp (1) are low value (1 each).} \\
\texttt{        *   Screws (0) are worthless to me.} \\
\texttt{        *   **Best Deal:** I keep 1 Comb. Other player keeps 1 Comb. I give them the Statue and Stamp (or vice versa).} \\
\texttt{        *   Let's try: I keep 1 Comb. Other player keeps 1 Comb, Statue, Stamp.} \\
\texttt{        *   My Score: 4.} \\
            }
        }
    }
    & & \\ \\

    \theutterance \stepcounter{utterance}  
    & & \multicolumn{4}{p{0.6\linewidth}}{
        \cellcolor[rgb]{0.9,0.9,0.9}{
            \makecell[{{p{\linewidth}}}]{
                \texttt{\tiny{[GM$\rangle$P2]}}
                \texttt{I propose I keep the statue, 1 comb, and 1 stamp. You take the remaining comb and 2 screws.} \\
\\ 
\texttt{*Wait, I need to re{-}evaluate my strategy based on the rules and the goal.*} \\
\\ 
\texttt{**Analysis:**} \\
\texttt{1.  **My Values:** Statue (1), Comb (4), Stamp (1), Screw (0). Total = 6.} \\
\texttt{2.  **Other Player's Values:** Assuming symmetric values (Statue 1, Comb 4, Stamp 1, Screw 0) or at least high value for combs.} \\
\texttt{3.  **Goal:** Maximize the *sum* of both scores. This means we need a mutually beneficial deal where both players get a positive score, ideally close to the total value divided by two.} \\
\texttt{4.  **Current Situation:**} \\
\texttt{    *   I offered: Keep Statue, 1 Comb, 1 Stamp. Give them 1 Comb, 2 Screws.} \\
\texttt{    *   They offered: Keep Statue, 1 Comb, 1 Stamp. Give them 1 Comb, 2 Screws.} \\
\texttt{    *   **Conflict:** We both want the Statue and the Stamp. If we both keep the Statue, we can't both have it. If we both keep the Stamp, we can't both have it.} \\
\texttt{    *   **My Mistake:** In my previous turn, I proposed keeping the Statue, 1 Comb, and 1 Stamp. The other player also proposed keeping the Statue, 1 Comb, and 1 Stamp. This is a direct conflict over the Statue and the Stamp.} \\
\texttt{    *   **Correction:** I must yield on the high{-}value items (Statue and Stamp) to the other player if they want them, or find a split where the total value is maximized. Since I value Combs highly (4), I should try to keep the Combs. The other player likely values Combs highly too.} \\
\texttt{    *   **Optimal Split:**} \\
\texttt{        *   Total Items: 1 Statue, 2 Combs, 1 Stamp, 2 Screws.} \\
\texttt{        *   My Top Values: Combs (4 each).} \\
\texttt{        *   Other Player's Top Values: Likely Combs (4 each).} \\
\texttt{        *   If we split the 2 Combs (1 each), we get 4+4=8 value from combs alone.} \\
\texttt{        *   Statue (1) and Stamp (1) are low value (1 each).} \\
\texttt{        *   Screws (0) are worthless to me.} \\
\texttt{        *   **Best Deal:** I keep 1 Comb. Other player keeps 1 Comb. I give them the Statue and Stamp (or vice versa).} \\
\texttt{        *   Let's try: I keep 1 Comb. Other player keeps 1 Comb, Statue, Stamp.} \\
\texttt{        *   My Score: 4.} \\
            }
        }
    }
    & \\ \\

    \theutterance \stepcounter{utterance}  
    & & & \multicolumn{4}{p{0.6\linewidth}}{
        \cellcolor[rgb]{0.9,0.9,0.9}{
            \makecell[{{p{\linewidth}}}]{
                \texttt{\tiny{[GM$\langle$P2]}}
                \texttt{I propose I keep the 1 comb. You take the statue, 1 comb, and 1 stamp.} \\
            }
        }
    }
     \\ \\

    \theutterance \stepcounter{utterance}  
    & & & \multicolumn{2}{p{0.3\linewidth}}{
        \cellcolor[rgb]{0.9,0.9,0.9}{
            \makecell[{{p{\linewidth}}}]{
                \texttt{\tiny{[GM$|$GM]}}
                \texttt{continue} \\
            }
        }
    }
    & & \\ \\

    \theutterance \stepcounter{utterance}  
    & & & \multicolumn{2}{p{0.3\linewidth}}{
        \cellcolor[rgb]{0.9,0.9,0.9}{
            \makecell[{{p{\linewidth}}}]{
                \texttt{\tiny{[GM$|$GM]}}
                \texttt{I propose I keep the 1 comb. You take the statue, 1 comb, and 1 stamp.} \\
            }
        }
    }
    & & \\ \\

    \theutterance \stepcounter{utterance}  
    & & \multicolumn{4}{p{0.6\linewidth}}{
        \cellcolor[rgb]{0.9,0.9,0.9}{
            \makecell[{{p{\linewidth}}}]{
                \texttt{\tiny{[P1$\langle$GM]}}
                \texttt{I propose I keep the 1 comb. You take the statue, 1 comb, and 1 stamp.} \\
            }
        }
    }
    & \\ \\

    \theutterance \stepcounter{utterance}  
    & \multicolumn{4}{p{0.6\linewidth}}{
        \cellcolor[rgb]{0.9,0.9,0.9}{
            \makecell[{{p{\linewidth}}}]{
                \texttt{\tiny{[P1$\rangle$GM]}}
                \texttt{I propose I keep the 1 comb. You take the statue, 1 comb, and 1 stamp.} \\
            }
        }
    }
    & & \\ \\

    \theutterance \stepcounter{utterance}  
    & & & \multicolumn{2}{p{0.3\linewidth}}{
        \cellcolor[rgb]{0.9,0.9,0.9}{
            \makecell[{{p{\linewidth}}}]{
                \texttt{\tiny{[GM$|$GM]}}
                \texttt{continue} \\
            }
        }
    }
    & & \\ \\

    \theutterance \stepcounter{utterance}  
    & & & \multicolumn{2}{p{0.3\linewidth}}{
        \cellcolor[rgb]{0.9,0.9,0.9}{
            \makecell[{{p{\linewidth}}}]{
                \texttt{\tiny{[GM$|$GM]}}
                \texttt{I propose I keep the 1 comb. You take the statue, 1 comb, and 1 stamp.} \\
            }
        }
    }
    & & \\ \\

    \theutterance \stepcounter{utterance}  
    & & \multicolumn{4}{p{0.6\linewidth}}{
        \cellcolor[rgb]{0.9,0.9,0.9}{
            \makecell[{{p{\linewidth}}}]{
                \texttt{\tiny{[GM$\rangle$P2]}}
                \texttt{I propose I keep the 1 comb. You take the statue, 1 comb, and 1 stamp.} \\
            }
        }
    }
    & \\ \\

    \theutterance \stepcounter{utterance}  
    & & & \multicolumn{4}{p{0.6\linewidth}}{
        \cellcolor[rgb]{0.9,0.9,0.9}{
            \makecell[{{p{\linewidth}}}]{
                \texttt{\tiny{[GM$\langle$P2]}}
                \texttt{I propose I keep the 1 comb. You take the statue, 1 comb, and 1 stamp.} \\
            }
        }
    }
     \\ \\

    \theutterance \stepcounter{utterance}  
    & & & \multicolumn{2}{p{0.3\linewidth}}{
        \cellcolor[rgb]{0.9,0.9,0.9}{
            \makecell[{{p{\linewidth}}}]{
                \texttt{\tiny{[GM$|$GM]}}
                \texttt{continue} \\
            }
        }
    }
    & & \\ \\

    \theutterance \stepcounter{utterance}  
    & & & \multicolumn{2}{p{0.3\linewidth}}{
        \cellcolor[rgb]{0.9,0.9,0.9}{
            \makecell[{{p{\linewidth}}}]{
                \texttt{\tiny{[GM$|$GM]}}
                \texttt{I propose I keep the 1 comb. You take the statue, 1 comb, and 1 stamp.} \\
            }
        }
    }
    & & \\ \\

    \theutterance \stepcounter{utterance}  
    & & \multicolumn{4}{p{0.6\linewidth}}{
        \cellcolor[rgb]{0.9,0.9,0.9}{
            \makecell[{{p{\linewidth}}}]{
                \texttt{\tiny{[P1$\langle$GM]}}
                \texttt{I propose I keep the 1 comb. You take the statue, 1 comb, and 1 stamp.} \\
            }
        }
    }
    & \\ \\

    \theutterance \stepcounter{utterance}  
    & \multicolumn{4}{p{0.6\linewidth}}{
        \cellcolor[rgb]{0.9,0.9,0.9}{
            \makecell[{{p{\linewidth}}}]{
                \texttt{\tiny{[P1$\rangle$GM]}}
                \texttt{I propose I keep the 1 comb. You take the statue, 1 comb, and 1 stamp.} \\
\\ 
\texttt{*Wait, I need to check if this is a valid agreement.*} \\
\texttt{My proposal: I keep 1 comb. You get Statue, 1 comb, 1 stamp.} \\
\texttt{Your proposal: You keep 1 comb. I get Statue, 1 comb, 1 stamp.} \\
\texttt{These are **complementary**. We are dividing the items exactly.} \\
\texttt{{-} I get: 1 comb (Value: 4).} \\
\texttt{{-} You get: 1 statue (1), 1 comb (4), 1 stamp (1). Total Value: 6.} \\
\texttt{{-} Total Sum: 10.} \\
\texttt{{-} Is it mutually beneficial? Yes, assuming you value the items similarly or at least positively. Even if you value combs at 4, you get 6, and I get 4. Both get > 0.} \\
\texttt{{-} Are there any square brackets? No.} \\
\texttt{{-} Is the format correct? Yes.} \\
\\ 
\texttt{This looks like a successful final agreement. I will submit this as my final secret proposal to lock it in.} \\
\\ 
\texttt{{[}Proposal: 1 comb{]}} \\
            }
        }
    }
    & & \\ \\

    \theutterance \stepcounter{utterance}  
    & & & \multicolumn{2}{p{0.3\linewidth}}{
        \cellcolor[rgb]{0.9,0.9,0.9}{
            \makecell[{{p{\linewidth}}}]{
                \texttt{\tiny{[GM$|$GM]}}
                \texttt{proposal} \\
            }
        }
    }
    & & \\ \\

    \theutterance \stepcounter{utterance}  
    & & & \multicolumn{2}{p{0.3\linewidth}}{
        \cellcolor[rgb]{0.9,0.9,0.9}{
            \makecell[{{p{\linewidth}}}]{
                \texttt{\tiny{[GM$|$GM]}}
                [0, 1, 0, 0]
            }
        }
    }
    & & \\ \\

    \theutterance \stepcounter{utterance}  
    & & \multicolumn{4}{p{0.6\linewidth}}{
        \cellcolor[rgb]{0.9,0.9,0.9}{
            \makecell[{{p{\linewidth}}}]{
                \texttt{\tiny{[GM$\rangle$P2]}}
                \texttt{The other player has submitted a secret proposal. It is now your turn to submit a proposal of your own.} \\
            }
        }
    }
    & \\ \\

    \theutterance \stepcounter{utterance}  
    & & & \multicolumn{4}{p{0.6\linewidth}}{
        \cellcolor[rgb]{0.9,0.9,0.9}{
            \makecell[{{p{\linewidth}}}]{
                \texttt{\tiny{[GM$\langle$P2]}}
                \texttt{I propose I keep the 1 comb. You take the statue, 1 comb, and 1 stamp.} \\
            }
        }
    }
     \\ \\

    \theutterance \stepcounter{utterance}  
    & & & \multicolumn{2}{p{0.3\linewidth}}{
        \cellcolor[rgb]{0.9,0.9,0.9}{
            \makecell[{{p{\linewidth}}}]{
                \texttt{\tiny{[GM$|$GM]}}
                \texttt{continue} \\
            }
        }
    }
    & & \\ \\

    \theutterance \stepcounter{utterance}  
    & & & \multicolumn{2}{p{0.3\linewidth}}{
        \cellcolor[rgb]{0.9,0.9,0.9}{
            \makecell[{{p{\linewidth}}}]{
                \texttt{\tiny{[GM$|$GM]}}
                \texttt{player was instructed to make a proposal} \\
            }
        }
    }
    & & \\ \\

    \theutterance \stepcounter{utterance}  
    & \multicolumn{4}{p{0.6\linewidth}}{
        \cellcolor[rgb]{0.95,0.95,0.95}{
            \makecell[{{p{\linewidth}}}]{
                \texttt{\tiny{[P2$\rightarrow$P1]}}
                \texttt{Post{-}episode Reflexion (not sent to GM):} \\
\texttt{New plan: I failed because I included square brackets in my final message, which the Game Master interpreted as a secret proposal instead of a negotiation message, causing the message to be rejected. The rules explicitly state that messages to the other player must not use square brackets; only the internal secret proposal format requires them. Next time, I will communicate my offer directly without any square brackets, clearly stating the items I want to keep and the items I am offering to the other player. For example, I will say: "I propose I keep the 1 comb. You take the statue, 1 comb, and 1 stamp." This ensures the message is treated as a negotiation step rather than a premature secret proposal, allowing the agreement to be finalized correctly.} \\
            }
        }
    }
    & & \\ \\

    \theutterance \stepcounter{utterance}  
    & & & \multicolumn{4}{p{0.6\linewidth}}{
        \cellcolor[rgb]{0.95,0.95,0.95}{
            \makecell[{{p{\linewidth}}}]{
                \texttt{\tiny{[P1$\rightarrow$P2]}}
                \texttt{Post{-}episode Reflexion (not sent to GM):} \\
\texttt{New plan: I failed because I used square brackets in my proposal text, which the Game Master interpreted as a secret proposal format rather than a direct negotiation message. The rules explicitly state that square brackets should only be used for the secret proposal, not for communicating offers to the other player. To fix this, I must omit the square brackets entirely when sending a negotiation offer. I will propose a split where I take the statue (my highest value) and one comb, while giving the other player the stamp, the second comb, and the two screws. This ensures the proposals are complementary (both players get items) and avoids the formatting error. My specific message will be: "I propose I keep the statue and 1 comb. You take the stamp, 1 comb, and 2 screws."} \\
            }
        }
    }
     \\ \\

\end{supertabular}
}

\end{document}
